\chapter{CONTENIDOS DEL TFG}
\thispagestyle{empty}

% =====================================================================
%  PENDIENTE DE REDACCION  --  Capitulo 2
%  Este capitulo contendra, en su version final, los siguientes
%  apartados (segun los indicativos de la plantilla de la Facultad):
%
%    \section{Marco Teorico}
%      - Redes Programables (SDN): arquitectura, plano de control y plano
%        de datos, controladores OpenFlow, Open vSwitch.
%      - Deteccion de intrusiones: firmas vs. anomaly detection.
%      - Machine Learning aplicado a seguridad de redes: LightGBM,
%        extraccion de caracteristicas, desbalance de clases, metricas.
%      - Plataformas de visualizaci\u00f3n y auditoria de eventos.
%
%    \section{Evaluacion de Riesgos}
%      - Alternativas consideradas (IDS con firmas, IPS comercial,
%        plataforma propia) y matriz de riesgos / ventajas.
%
%    \section{Resultados}
%      - Escenario de pruebas, metricas (detecci\u00f3n, falsos positivos,
%        tiempo de mitigaci\u00f3n) y an\u00e1lisis de resultados.
%
%    \section{Conclusiones y Lineas Futuras}
%      - Ver cap04.tex.
%
%  While  se escribe este capitulo, puede comentarse la inclusion desde
%  ReportePI.tex con \includeonly{cap01} para acelerar la compilacion.
% =====================================================================

\section{Marco Te\'orico}
% --- PENDIENTE: redactar el marco teorico del proyecto ---

\section{Evaluaci\'on de Riesgos}
% --- PENDIENTE: comparar alternativas tecnologicas y evaluar riesgos ---

\section{Resultados}
% --- PENDIENTE: presentar los resultados experimentales obtenidos ---

\section{Conclusiones}
% --- PENDIENTE: ver cap04.tex (CONCLUSIONES Y LINEAS FUTURAS) ---

\section{L\'ineas Futuras}
% --- PENDIENTE: plantear las lineas futuras de trabajo ---
